\section{Results}
\label{sec:results}

\subsection{The app as shipped}
\label{sec:res-shipped}

Given the system prompt alone, Llama~3.1~8B returned an empty answer for all \ShippedEmpty{}
questions (Table~\ref{tab:answers}). In the app, this is what the first question of every
conversation received: a blank reply, with nothing to show that the question and the document
had not been sent.

\begin{table}[t]
\centering
\caption{Answers on the QASPER development set, judged by Qwen3-32B (Wilson 95\% intervals).
\emph{Shipped}: the system prompt alone. \emph{Fixed}: the fixed app's prompt with the four
retrieved chunks; the three questions on the paper whose PDF could not be downloaded are left
out. \emph{Oracle}: the same prompt with the annotators' text evidence, for the questions that
have some. Answer F1 is the official QASPER metric, which compares tokens with short reference
answers.}
\label{tab:answers}
\small
\begin{tabular}{lrccc}
\toprule
Condition & Questions & Correct (\%) & Correct or partial (\%) & Answer F1 \\
\midrule
\rows{generated/answers_rows}
\bottomrule
\end{tabular}
\end{table}

\subsection{Retrieval limits the fixed app}
\label{sec:res-fixed}

With the bug fixed, \FixedCorrect{} of the \FixedN{} answers were judged correct and
\FixedPartial{} correct or partially correct. Official Answer F1 was \FixedFOne{}; it is low for
every condition because the app answers in full sentences and the metric counts the words that
are not in a short reference answer, so we use it only alongside the judged rates.

Replacing the retrieved chunks with the annotators' evidence paragraphs, on the \PairedOracleN{}
questions that have text evidence and were answered in both conditions, raised the rate from
\PairedFixedOnOracleQuestions{} to \OracleCorrect{} (\PairedOracleGained{} questions gained,
\PairedOracleLost{} lost; exact McNemar $p =$ \PairedOracleP{}). The model and prompt were the
same, so the difference is due to what the context held.

Two things limit the retrieved context. A chunk of the downloaded PDF held the gold evidence for
\Findable{} of the \TextEvidenceQuestions{} questions with text evidence; the rest were lost to
PDF extraction or to differences between the arXiv version and the text QASPER was built on.
Of the \FindableK{} findable questions, the four retrieved chunks held evidence for \HitFour{},
and the top ten for \HitTen{}. The fixed app was judged correct for \CorrectWhenHit{} of the
\CorrectWhenHitN{} questions whose four chunks held evidence and for \CorrectWhenMiss{} of the
\CorrectWhenMissN{} whose chunks did not. That second figure is not small: some questions can
be answered from text near the evidence, and some answers were judged correct for saying the
paper did not answer the question. The comparison is observational, since questions whose
evidence is easy to retrieve may also be easier to answer; the oracle comparison above is the
controlled one.

\subsection{Retrieval settings at equal prompt length}
\label{sec:res-sweep}

With the prompt held at about 4,000 characters, the app's set-up (1,000-character chunks, top 4,
MiniLM) put a gold evidence passage in the context for \SweepBaseBudget{} of the
\SweepQuestions{} questions (Table~S\result{n}). Halving the chunk size and taking twice as many
chunks, with the same encoder, raised this to \SweepSmallBudget{} (\SweepSmallGained{} questions
gained, \SweepSmallLost{} lost; Holm $p =$ \SweepSmallP{}). Reranking the top 20 of those
500-character chunks with the cross-encoder raised it to \SweepSmallRerankBudget{}
(\SweepSmallRerankGained{} gained, \SweepSmallRerankLost{} lost against the app's set-up; Holm
$p =$ \SweepSmallRerankP{}). The best of the \SweepSetups{} set-ups, \SweepBestName{}, reached
\SweepBestBudget{}, inside the interval of the simpler reranked set-up; with reranking, the
choice of encoder and the addition of BM25 changed little. BM25 alone was the worst retriever at
every chunk size.

These are rates of evidence reaching the model, and the matching rule may favour small chunks
(a 500-character chunk inside a long evidence paragraph can pass the overlap test while holding
only part of it). We therefore tested the two leading set-ups on answers.

\subsection{The leading set-ups on answers}
\label{sec:res-candidates}

\result{Fixed app (c1000, top 4) vs c500 top 8 vs c500 top 8 reranked: judged correct with
intervals, paired McNemar against the fixed app, official F1, mean prompt length, seconds per
question.}

\subsection{Several papers}
\label{sec:res-multidoc}

Other papers crowded out the one a question was about (Table~\ref{tab:multidoc}). For the
\MultidocQuestions{} questions with findable evidence, the top four chunks held evidence for
\MultidocZeroSingle{} with the paper alone. With nine other papers loaded, only
\MultidocNineSingleOwn{} of the four chunks came from the right paper on average, and the rate
fell to \MultidocNineSingle{} (\MultidocNineSingleLost{} questions lost, none gained). The rule
the app uses for several files, which adds one chunk from each unrepresented file, recovered
part of this mainly by widening the context to eight chunks (\MultidocNineMulti{}).
Restricting the search to the question's paper restored the single-paper rate exactly, as it
must, since it searches the same chunks with the same embeddings.

\begin{table}[t]
\centering
\caption{Evidence retrieval when the question's paper is loaded with $N$ other QASPER
development papers (\MultidocQuestions{} questions with findable evidence; Wilson 95\%
intervals). \emph{single}: top 4 chunks. \emph{multi}: top 4, then the best chunk of each other
file, up to 8. \emph{selected}: top 4 of the question's paper only. Lost / gained: questions
whose evidence left or entered the context compared with the paper alone.}
\label{tab:multidoc}
\small
\begin{tabular}{rlcccc}
\toprule
$N$ & Rule & Evidence in context (\%) & Chunks & From right paper & Lost / gained \\
\midrule
\rows{generated/multidoc_rows}
\bottomrule
\end{tabular}
\end{table}

\subsection{Agreement of the judge with hand labels}
\label{sec:res-judge}

On the \JudgeCheckN{} sampled answers, Claude's grades and the judge's agreed exactly for
\JudgeCheckExact{} (Cohen's $\kappa =$ \JudgeCheckKappa{}) and on correct versus not correct for
\JudgeCheckBinary{} ($\kappa =$ \JudgeCheckKappaBinary{}). The judge called \JudgeCheckJudgeCorrect{}
of them correct and Claude \JudgeCheckClaudeCorrect{}. Its errors ran both ways. In two clear
cases it credited answers that did not answer (``the paper does not say'' accepted as ``no'', and
a hedge accepted as ``yes''), and in four it called a complete answer partial because it added
or left out detail (every disagreement is listed in the SI). This is agreement between two models, not
accuracy against people; \result{agreement of each with human labels on the same answers}.
